\documentclass[12pt,a4paper]{article}
\usepackage[utf8]{inputenc} %codificacao de entrada
\usepackage[T1]{fontenc} %codificacao de saida
\usepackage[brazil]{babel}

\usepackage{titling}		% arruma pre titulo
\usepackage{makeidx}        % Cria o indice
\usepackage{hyperref}       % Controla a formacao do indice
\usepackage{lastpage}       % Usado pela Ficha catalografica
\usepackage{indentfirst}    % Indenta o primeiro paragrafo de cada secao.
\usepackage{xcolor}          % Controle das cores
\usepackage{graphicx}       % Inclusao de graficos
\usepackage{placeins}
\usepackage{subfig}
\usepackage{amsmath}        % pacote matemático
\usepackage{listings}
\usepackage{longtable}

\usepackage[hmargin=2cm,vmargin=3.5cm,bmargin=2cm]{geometry}

\pretitle{%
  \begin{center}
  \LARGE
  \includegraphics[width=3.5cm,height=4cm]{imgs/logo.png}\\[1.5cm]\textbf{Universidade do Estado do Rio de Janeiro}\\ Instituto Politécnico\\Engenharia de Computação\\[2cm]
}
\posttitle{\normalsize \end{center}}

\title{\textbf{\underline{Trabalho}\\[0.5cm]
Processos Estocásticos}\\[0.3cm]\large{Professor Angelo}\\[1.5cm]}
\author{
	Ana Carolina Castro\\
	\texttt{anacarolinacastro@aol.com}
	\and
	Ciro Chang\\
	\texttt{cirochang@live.com}
	\\[1.5cm]
}
\date{Nova Friburgo\\31 de Outubro de 2015}

    
\begin{document}
\maketitle
\thispagestyle{empty}

\newpage
\tableofcontents

\newpage
\section{Introdução}
Esse trabalho tem como objetivo analisar os resultados de séries aleatórias. Usaremos quatro séries, onde
a primeira e a segunda serão compostas por valores totalmente aleatórios entre $-1$ e $1$. A terceira será uma cópia parcial da primeira série, isto é, para cada valor da primeira séries iremos tirar na sorte se o valor será copiado ou um novo valor aleatório. E a quarta série será dada pela função $sen(x)$, para cada valor da primeira série.\\[1cm]

\section{O Caminhante Bebado}
\subsection{Variáncia}
\subsection{1 Dimensão}
\subsection{2 Dimensões}
hehehe\\[1.5cm]

\section{Gota}
A gota cai, a rua inunda.
\subsection{Simulação}
\subsection{Experimento}

\end{document}




